\chapter{Glossary of recurring structures}
\label{chap:glossary}

\begin{description}[style=nextline]
  \item[Aguilar--Balslev--Combes theorem (ABC theorem).]
  Under dilation-analytic and relative-compactness hypotheses, complex
  dilation leaves bound eigenvalues fixed, rotates the essential spectrum,
  and exposes resonance poles as discrete eigenvalues independent of the
  dilation angle inside their exposure wedge.  The theorem concerns an
  analytic operator family and is stronger than the visual stabilization of
  a finite matrix spectrum.

  \item[Adjoint operator.]
  For a densely defined operator $A$, the domain of $A^*$ consists of vectors
  for which $\psi\mapsto\dualpair{\phi}{A\psi}$ is continuous in the
  Hilbert norm.  The domain is part of the definition; formal integration by
  parts alone does not prove that $A=A^*$.
  \item[Analytic dilation.] Continuation of a real coordinate dilation to a
  complex parameter in a sector where the Hamiltonian is an analytic family.
  It is theorem-controlled and is not synonymous with adding an absorber.

  \item[Analytic family of operators.]
  A parameter-dependent family whose resolvent, quadratic form, or action on
  a common domain is analytic in the sense specified by Kato.  Stating the
  type and domain prevents a moving numerical matrix from being mistaken for
  analytic deformation of an unbounded Hamiltonian.

  \item[Analytic Fredholm theorem.]
  For an analytic compact-operator family $K(z)$, either
  $\identity-K(z)$ is nowhere invertible or its inverse is meromorphic, with
  finite-rank principal parts at isolated poles.  This is the operator-level
  bridge from an integral equation to a resonance pole.

  \item[Antiresonance.]
  The pole or connection zero paired with a resonance by the relevant
  reality, time-reversal, or sheet symmetry.  It is not a second decaying
  state on the same physical boundary-value problem.

  \item[Asymptotic completeness.]
  The assertion that the ranges of the incoming and outgoing wave operators,
  together with bound states, exhaust the physical Hilbert space.  Existence
  of M{\o}ller operators is weaker than completeness.

  \item[Berggren completeness relation.]
  A contour-deformed resolution combining discrete bound and resonant poles
  with a nonresonant continuum integral.  Moving the contour can change which
  poles appear explicitly while leaving a complete observable unchanged.

  \item[Biorthogonal system.]
  Right and left vectors $\psi_n$ and $\widetilde\psi_n$ satisfying
  $\braket{\widetilde\psi_m|\psi_n}=\delta_{mn}$ where diagonalization is
  possible.  Near an exceptional point the normalization becomes singular
  and must be replaced by Jordan-chain data.

  \item[Birman--Schwinger operator.]
  A compactly sandwiched resolvent, for example
  $K(z)=|V|^{1/2}(H_0-z)^{-1}\operatorname{sgn}(V)|V|^{1/2}$.
  Under the stated domain and compactness hypotheses, $z$ is an eigenvalue
  or continued pole exactly when the corresponding $K(z)$ has eigenvalue
  one.

  \item[Borel transform and lateral sum.]
  The Borel transform divides the coefficient of $\hbar^n$ by a factorial so
  that a divergent formal series may acquire a finite convergence radius.
  Laplace integration just above or below a singular Borel ray defines the
  two lateral sums; their difference is a Stokes discontinuity.

  \item[Breit--Wigner approximation.]
  The Lorentzian real-axis form obtained from an isolated simple pole only
  after the background, residue, phase space, and self-energy are treated as
  slowly varying.  Pole energy and fitted peak parameters need not agree for
  a broad or threshold-adjacent resonance.

  \item[Branch cut and sheet.]
  A cut is a convention used to make a multivalued analytic function
  single-valued on one chart; a sheet is the corresponding branch of its
  Riemann surface.  Resonance statements must specify continuation relative
  to thresholds, not merely say ``the second sheet.''

  \item[Channel.]
  A chosen asymptotic internal state, relative orbital motion, and set of
  conserved quantum numbers.  An open channel carries flux at the selected
  energy; a closed channel decays asymptotically but can still generate a real
  polarization potential.

  \item[Continuum-discretized coupled channels (CDCC).]
  A projection method that replaces selected projectile-continuum sectors by
  normalized bins or pseudostates and solves the resulting coupled radial
  equations.  Convergence is a statement about observables under enlargement
  of the model-space projection, not about individual positive-energy
  eigenvalues.

  \item[Continuum level density (CLD).]
  The imaginary part of a reference-subtracted resolvent trace.  Under
  trace-class hypotheses it is the derivative of a scattering spectral shift
  and includes both pole and nonresonant continuum information.

  \item[Complex-scaled Lippmann--Schwinger method (CSLS).]
  A representation of a scattering resolvent in a complex-scaled basis,
  commonly used to smooth CDCC pseudostate amplitudes and resolve breakup
  spectra.  Angle independence is recovered only after a sufficiently
  complete basis and consistent inverse transformation.

  \item[Complex scaling.]
  Analytic continuation of a dilation, commonly $r\mapsto r\rme^{\rmi\theta}$,
  which converts an exposed outgoing resonance into an $L^2$ eigenproblem and
  rotates continuum cuts.  Exterior and smooth variants modify the map while
  preserving the same analytic aim.

  \item[Connection coefficient.] A scalar minor or matrix block relating
  local solution bases in distinct asymptotic sectors.  Its prohibited-
  incoming component is denoted $\Cincoming$.

  \item[Continuous spectrum.]
  The spectral part in which $H-E$ fails to have a bounded inverse although
  $E$ is not a normalizable eigenvalue.  Generalized energy eigenvectors are
  distributions or fiber labels, not additional vectors of $\Hil$.

  \item[Continuum pseudostate.] A positive-energy $L^2$ eigenvector of a
  finite-basis diagonalization used as a quadrature state.  Its individual
  energy is normally representation dependent.

  \item[$c$ product.]
  The bilinear pairing natural for a complex-symmetric matrix or a suitably
  scaled differential operator.  It transposes rather than Hermitian
  conjugates the coefficient vector and must not be used when the operator's
  left problem has a different structure.

  \item[Direct integral.] A measure-theoretic continuum of Hilbert-space
  fibers, used here to represent energy and channel multiplicity.

  \item[Dollard comparison dynamics.]
  A long-range modifier of free motion that includes the logarithmic Coulomb
  phase before taking the scattering limit.  It repairs the failure of the
  ordinary free M{\o}ller operator for Coulomb interactions.

  \item[Dressing sector.]
  An infrared representation in which charged asymptotic states carry a
  correlated cloud of soft gauge quanta.  The required profile generally has
  infinite norm in the naive Fock one-particle space, so changing charge or
  velocity can move between inequivalent infrared sectors.

  \item[Effect and POVM.]
  An effect is an operator $0\le F\le\identity$ representing one measurement
  outcome.  A positive-operator-valued measure assigns effects to measurable
  outcome sets and includes finite-resolution measurements that are not
  projection-valued.

  \item[Essential spectrum.]
  The part of the spectrum stable under the appropriate compact
  perturbations.  Complex scaling rotates its threshold rays; finite-basis
  points sampling those rays are pseudostates rather than isolated spectrum
  of the exact operator.

  \item[Exceptional point (EP).]
  A parameter value at which eigenvalues and eigenvectors coalesce so that
  algebraic multiplicity exceeds geometric multiplicity.  In the connection
  formulation it is a multiple zero supplemented by a defectiveness test.

  \item[Feshbach operator.]
  The energy-dependent operator on a retained $P$ space obtained by solving
  the $Q$ equation with a specified resolvent boundary value.  Its derivative
  enters pole normalization because the eliminated sector remains encoded in
  the self-energy.

  \item[Fano profile.]
  The asymmetric intensity $|A_0|^2(\epsilon+q)^2/(1+\epsilon^2)$ produced
  by coherent interference between a pole amplitude and a smooth
  background.  Its maximum and zero are reaction dependent.

  \item[Fiber Hilbert space.]
  The multiplicity space $\Hil_E$ attached to almost every energy in a direct
  integral.  Channel labels live in the fiber; measurable wave packets are
  square-integrable sections $E\mapsto\psi(E)$.

  \item[Gamow--Siegert state.] A solution satisfying outgoing boundary
  conditions at complex energy.  It is generally not in the original
  Hilbert space.

  \item[Gelfand--Naimark--Segal construction (GNS).]
  The construction of a Hilbert-space representation from a positive linear
  functional on a $C^*$ algebra.  It makes explicit that a state is primary
  and a chosen Hilbert-space representation can depend on its physical
  sector.

  \item[Infrared divergence.]
  Failure of an integral at vanishing momentum or energy.  In long-range
  scattering it can signal that the assumed asymptotic state space is wrong,
  rather than merely that a finite parameter needs renormalization.

  \item[Hilbert--Schmidt operator.]
  A compact operator whose squared singular values are summable.  An
  integral kernel $K(x,y)$ is Hilbert--Schmidt when
  $\int|K(x,y)|^2\rmd x\rmd y<\infty$; then finite-rank kernel
  approximations converge in Hilbert--Schmidt, hence operator, norm.

  \item[Jost function.]
  A Wronskian coefficient between a solution regular at an interior endpoint
  and an asymptotic incoming/outgoing basis.  Under suitable short-range and
  analyticity assumptions, its zeros identify bound, virtual, resonant, or
  antiresonant poles according to their sheet and quadrant.

  \item[Jordan chain.]
  Vectors satisfying $(H-E_*)\psi_0=0$ and
  $(H-E_*)\psi_{j+1}=\psi_j$.  A chain of length two produces both simple and
  double terms in the resolvent principal part.

  \item[Langer correction.]
  The $l(l+1)\mapsto(l+1/2)^2$ term obtained when the radial equation is
  transformed logarithmically and its wave function is treated as a
  half-density.  It records origin monodromy rather than an empirical WKB
  patch.

  \item[M{\o}ller or wave operator.]
  The strong limit comparing full and reference time evolution at early or
  late times.  Its definition includes the comparison Hamiltonian; for
  Coulomb or QED asymptotics, free Fock or free-particle comparison may be
  inadequate.

  \item[Monodromy.]
  The linear transformation of a local solution basis after analytic
  continuation around a closed path.  Turning points, regular singular
  points, and coordinate Jacobians contribute distinguishable factors.

  \item[Nuclear optical potential.]
  A generally complex, energy-dependent effective interaction for an elastic
  or selected channel.  Its absorptive part represents flux into eliminated
  channels and should not be confused with fundamental nonunitarity of the
  full many-body Hamiltonian.

  \item[Projection-valued measure (PVM).]
  The countably additive family of orthogonal projections supplied by the
  spectral theorem for a self-adjoint operator.  Matrix elements of the PVM
  define scalar spectral measures and direct-integral decompositions.

  \item[Quasinormal mode (QNM).]
  A mode satisfying ingoing conditions at a future black-hole horizon and
  outgoing or appropriate cosmological/AdS conditions at the other endpoint.
  With $\rme^{-\rmi\omega t}$, a decaying mode has $\im\omega<0$.

  \item[Resolvent.]
  $\Resolvent(z)=(z-H)^{-1}$ on the resolvent set.  Its boundary values encode the
  spectral measure; matrix elements or weighted versions may admit analytic
  continuation through a continuous-spectrum cut even though the bounded
  Hilbert-space operator does not.

  \item[Resonance.] A pole of an analytically continued resolvent or
  scattering matrix, equivalently a zero of the appropriate Jost or incoming
  connection function under the stated hypotheses.

  \item[Residue.]
  The coefficient of $(z-z_R)^{-1}$ in a Laurent expansion.  It contains the
  normalization-invariant product of left and right pole data and is required
  to predict response strength; a pole position alone is incomplete.

  \item[Representation datum.] A contour, scaling angle, test topology,
  basis, bin width, or infrared sector chosen to realize an invariant analytic
  quantity.

  \item[Riesz projection.]
  For an isolated spectral cluster inside a contour $\Gamma$, the operator
  $(2\pi\rmi)^{-1}\oint_\Gamma(z-H)^{-1}\rmd z$.  It is invariant under
  basis changes and treats a fragmented or nearly degenerate resonance as a
  subspace rather than as a preferred list of eigenvectors.

  \item[Rigged Hilbert space.]
  A continuous dense embedding $\PhiSpace\subset\Hil\subset\PhiSpace^\times$
  in which generalized eigenvectors act as continuous functionals on a test
  space.  The topology of $\PhiSpace$ determines which Gamow and scattering
  functionals exist.

  \item[$R$-matrix.]
  A boundary relation obtained by solving an interior problem at a chosen
  channel radius and matching it to known exterior functions.  Physical
  observables should become insensitive to the channel radius and boundary
  parameter in a complete expansion.

  \item[Self-adjoint extension.]
  A choice of operator domain, often encoded by boundary data, that makes a
  symmetric differential operator equal to its adjoint.  Different
  extensions can represent different physical contact or endpoint data.

  \item[Spectral measure.]
  For a vector $\psi$ and PVM $P(\cdot)$, the positive measure
  $\mu_\psi(B)=\dualpair{\psi}{P(B)\psi}$.  Survival amplitudes are Fourier
  transforms of such measures; exponential decay is therefore an
  intermediate-time approximation when a threshold is present.

  \item[Spectral shift.]
  The reference-subtracted change in spectral counting associated with a
  pair $(H,H_0)$.  Its derivative is a continuum level density and its
  exponential is related to $\det S$ by the Birman--Krein formula under the
  appropriate hypotheses.

  \item[Sturm--Liouville operator.]
  The differential expression
  $w^{-1}[-(pu')'+qu]$ together with its weighted Hilbert space, maximal
  domain, and endpoint boundary conditions.  The expression alone does not
  define a self-adjoint operator.

  \item[Stokes curve and Stokes matrix.]
  A Stokes curve is a locus where the phase relation between local WKB
  branches changes dominance.  Crossing it changes the resummed basis by a
  triangular Stokes matrix; the path-ordered product with propagation factors
  is the global connection matrix.

  \item[Threshold.]
  The endpoint of an asymptotic channel continuum and hence a branch point of
  its resolvent.  Threshold powers can dominate late-time decay and can
  compete with the square-root branching of a nearby exceptional point.

  \item[Topology.]
  A specification of which sets are open, equivalently which convergence and
  continuity statements are meaningful.  Norm, strong, weak, weak-star, and
  test-function topologies answer different physical questions and must not
  be silently interchanged.

  \item[\vN\ algebra.]
  A unital $*$ algebra of bounded operators closed in the weak-operator
  topology, equivalently equal to its bicommutant.  Abelian spectral algebras
  are multiplication algebras; their direct-integral fibers expose channel
  multiplicity and superselection structure.

  \item[Voros symbol.] The Borel-resummed exponential of an exact-WKB period
  integral on the spectral curve.

  \item[Wronskian.]
  The determinant of two local solutions and their derivatives.  It is
  constant for a second-order equation without a first-derivative term, fixes
  the Green-kernel denominator, and provides a sensitive determinant check
  for transfer and Stokes matrices.

  \item[Weyl--Titchmarsh function.]
  The coefficient $m(z)$ selecting the solution square integrable at a
  limit-point endpoint.  Its boundary value determines the spectral measure,
  while its analytically continued poles encode the same denominator as the
  Jost or Green-function construction.

  \item[Zel'dovich regularization.] Gaussian damping of a divergent outgoing
  pairing followed by analytic continuation in the regulator.
\end{description}
